\chapter{Modelling the Photovoltaic Array}
\label{ch:pvmodel}

\section{What the model has to be good at}

For an MPPT study the model must do more than reproduce the three datasheet
points. The controller lives on the \emph{curvature} of the characteristic: the
plant gain is set by the dynamic resistance $r_d$, which is a derivative, and
the MPPT layer operates on $\mathrm{d}P/\mathrm{d}V$, which is another
derivative. A model that passes through $\Isc$, $\Voc$ and the MPP but has the
wrong shape between them would give plausible-looking power numbers and
completely wrong control conclusions.

Three requirements follow:

\begin{enumerate}[leftmargin=1.6em]
  \item The five parameters must be identified from constraints that pin the
        \emph{slope} as well as the position of the curve.
  \item Derivatives must be computed analytically, not by finite differences of
        a numerically solved current.
  \item The maximum power point used as ground truth must be exact to solver
        tolerance, since it forms the denominator of the efficiency metric.
\end{enumerate}

\section{Parameter identification}

\subsection{The five constraints}

There are five unknowns, $\{I_{ph}, I_0, R_s, R_{sh}, a\}$. Five independent
conditions are imposed, all evaluated at STC:

\begin{align}
\text{short circuit:} \quad & I(0) = \Isc \label{eq:c1}\\
\text{open circuit:} \quad & I(\Voc) = 0 \label{eq:c2}\\
\text{maximum power point:} \quad & I(\Vmp) = \Imp \label{eq:c3}\\
\text{stationarity at the MPP:} \quad & \Imp + \Vmp \left.\frac{\mathrm{d}I}{\mathrm{d}V}\right|_{\Vmp} = 0 \label{eq:c4}\\
\text{temperature coefficient:} \quad & \left.\frac{\mathrm{d}\Voc}{\mathrm{d}T}\right|_{T_{ref}} = \beta_{Voc} \label{eq:c5}
\end{align}

Conditions~\eqref{eq:c1} to~\eqref{eq:c3} place the curve through three known
points. Condition~\eqref{eq:c4} is the crucial one: it forces the model to have
the correct \emph{slope} at the MPP, which is what fixes the curvature and hence
$r_d$. Condition~\eqref{eq:c5} closes the system using the published temperature
coefficient, which is the De Soto closure and uses information that simpler
identification schemes discard.

\subsection{Why not fix the ideality factor}

A common alternative, due to Villalva et al.~\cite{villalva}, fixes $a$ at an
assumed value and iterates on $R_s$ until the modelled peak power matches the
datasheet. That method is robust and widely used, but it leaves one parameter to
the modeller's judgement and does not use $\beta_{Voc}$ at all. Solving all five
equations simultaneously removes that arbitrariness, and the resulting ideality
factor becomes an \emph{output} that can be sanity-checked against physics: a
value near unity is expected for a good crystalline silicon cell.

\subsection{Numerical formulation}

The system is solved with a trust-region method on the residual vector. Three
implementation choices make it reliable:

\begin{itemize}[leftmargin=1.4em]
  \item \textbf{Log parameterisation.} The solver works with
        $\left[I_{ph},\ \ln I_0,\ \ln R_s,\ \ln R_{sh},\ n\right]$, which keeps
        every physical quantity positive by construction and conditions the
        problem, since $I_0$ spans ten orders of magnitude.
  \item \textbf{Residual scaling.} The four current residuals are divided by
        $\Isc$ and the temperature residual by $|\beta_{Voc}|$, so all five
        entries are of comparable size and the trust region is meaningful.
  \item \textbf{Multi-start.} The solve is repeated from six ideality-factor
        seeds between \num{0.8} and \num{2.0}, and the best converged result is
        kept, so the answer does not depend on a lucky initial guess.
\end{itemize}

\subsection{Results}

Table~\ref{tab:pvfit} shows the identified parameters and
Table~\ref{tab:pvvalid} the resulting agreement with the datasheet. All five
constraints are satisfied to machine precision, and the extracted ideality
factor of \num{1.0036} is physically sensible for a multicrystalline silicon
cell.

\begin{table}[htbp]
\centering
\caption{Identified reference parameters for the Kyocera KC200GT, module level, at standard test conditions.}
\label{tab:pvfit}
\begin{tabular}{@{}llll@{}}
\toprule
\textbf{Parameter} & \textbf{Symbol} & \textbf{Value} & \textbf{Unit} \\
\midrule
Photocurrent            & $I_{ph,ref}$  & \num{8.22713}    & \si{\ampere} \\
Saturation current      & $I_{0,ref}$   & \num{4.3964e-10} & \si{\ampere} \\
Series resistance       & $R_{s,ref}$   & \num{0.33501}    & \si{\ohm} \\
Shunt resistance        & $R_{sh,ref}$  & \num{160.6015}   & \si{\ohm} \\
Modified ideality factor& $a_{ref}$     & \num{1.39246}    & \si{\volt} \\
Diode ideality factor   & $n$           & \num{1.0036}     & -- \\
\bottomrule
\end{tabular}
\end{table}

\begin{table}[htbp]
\centering
\caption{Model against datasheet at standard test conditions. The residual is at the level of double-precision round-off.}
\label{tab:pvvalid}
\begin{tabular}{@{}lrrr@{}}
\toprule
\textbf{Quantity} & \textbf{Datasheet} & \textbf{Model} & \textbf{Error (\si{\percent})} \\
\midrule
$\Voc$ (\si{\volt})            & 32.90    & 32.90    & \num{0} \\
$\Isc$ (\si{\ampere})          & 8.21     & 8.21     & \num{0} \\
$\Vmp$ (\si{\volt})            & 26.30    & 26.30    & \num{-1.4e-14} \\
$\Imp$ (\si{\ampere})          & 7.61     & 7.61     & \num{0} \\
$\Pmp$ (\si{\watt})            & 200.14   & 200.14   & \num{0} \\
$\beta_{Voc}$ (\si{\volt\per\kelvin}) & -0.12301 & -0.12301 & \num{1.1e-14} \\
\bottomrule
\end{tabular}
\end{table}

The same procedure applied to a SunPower SPR-305E-WHT-D, a 96-cell
monocrystalline module of quite different construction, converges equally well
with a worst error of \SI{7e-11}{\percent}, which indicates that the method is
not tuned to one datasheet.

\section{Solving the implicit equation}
\label{sec:solver}

Because~\eqref{eq:singlediode} cannot be inverted in closed form, terminal
current is obtained by Newton iteration on the residual

\begin{equation}
f(I) = I_{ph} - I_0\left[\exp\!\left(\tfrac{V + I R_s}{a}\right) - 1\right] - \frac{V + I R_s}{R_{sh}} - I .
\end{equation}

Three properties make this safe. First, $f$ is strictly decreasing in $I$, so
the root is unique. Second, starting from $I = 0$ the residual is positive
whenever $V < \Voc$, and for a concave decreasing function Newton's method
converges monotonically from that side. Third, the exponent is clamped and a
backtracking line search halves the step whenever the residual magnitude would
grow, which covers the transient excursions beyond $\Voc$ that a controller can
briefly demand.

In simulation the previous sample's current is supplied as the initial guess.
This warm start reduces the iteration count to two or three and is what makes
second-scale closed-loop simulation practical.

\section{Derived quantities}

The slope is evaluated from the analytic expression~\eqref{eq:didv} rather than
by differencing. Verification against a central difference over the range
$0.05\,\Voc$ to $0.97\,\Voc$ agrees to a relative error below \num{1e-5}, which
is the accuracy floor of the difference itself, not of the analytic formula.

The maximum power point is found by scanning \num{400} points to bracket the
peak and then root-solving the analytic $\mathrm{d}P/\mathrm{d}V$ inside that
bracket. The stationarity residual at the returned point is below
\SI{1e-6}{\watt\per\volt}.

\section{Behaviour of the model}

\begin{figure}[htbp]
\centering
\includegraphics[width=\textwidth]{s01_pv_characteristics.png}
\caption[Photovoltaic array characteristics]{Characteristics of the 6 series by 2 parallel KC200GT array. Top left, I--V curves at \SI{25}{\celsius}; top right, the corresponding P--V curves with a single interior maximum; bottom left, the effect of cell temperature at \SI{1000}{\watt\per\square\metre}; bottom right, dynamic resistance $r_d$ on a logarithmic axis, which is the quantity that sets the plant gain.}
\label{fig:pvchar}
\end{figure}

Figure~\ref{fig:pvchar} shows the identified model in use.
Table~\ref{tab:pvops} lists the array operating points.

\begin{table}[htbp]
\centering
\caption{Array operating points at \SI{25}{\celsius}, 6 series by 2 parallel.}
\label{tab:pvops}
\begin{tabular}{@{}rrrrrrrr@{}}
\toprule
$G$ & $\Vmp$ & $\Imp$ & $\Pmp$ & $\Voc$ & $\Isc$ & $FF$ & $r_d$ at MPP \\
(\si{\watt\per\square\metre}) & (\si{\volt}) & (\si{\ampere}) & (\si{\watt}) & (\si{\volt}) & (\si{\ampere}) & & (\si{\ohm}) \\
\midrule
200  & 156.02 & 3.061  & 477.6  & 183.97 & 3.290  & 0.789 & 50.97 \\
400  & 158.75 & 6.117  & 971.1  & 189.75 & 6.576  & 0.778 & 25.95 \\
600  & 159.21 & 9.163  & 1458.8 & 193.14 & 9.860  & 0.766 & 17.38 \\
800  & 158.75 & 12.198 & 1936.4 & 195.54 & 13.141 & 0.754 & 13.02 \\
1000 & 157.80 & 15.220 & 2401.7 & 197.40 & 16.420 & 0.741 & 10.37 \\
\bottomrule
\end{tabular}
\end{table}

Three observations from these results shape everything that follows.

\subsection{The MPP voltage is remarkably stable against irradiance}

Across a fivefold change in irradiance, $\Vmp$ moves only from
\SI{156.0}{\volt} to \SI{159.2}{\volt}, about \SI{2}{\percent}. This is why
regulating PV \emph{voltage} is a sensible inner objective: the target barely
moves, so the outer MPPT layer only has to make small corrections. It is also
the basis of the crude ``fractional open-circuit voltage'' MPPT methods, which
simply hold $\vpv \approx 0.78\,\Voc$.

\subsection{Temperature moves the MPP far more than irradiance does}

The bottom-left panel of Figure~\ref{fig:pvchar} shows $\Vmp$ falling
substantially as cell temperature rises, because $\Voc$ falls at
\SI{-0.123}{\volt\per\kelvin} per module, that is \SI{-0.738}{\volt\per\kelvin}
for the six-module string. Over a realistic \SI{40}{\kelvin} swing that is
nearly \SI{30}{\volt}. Temperature, not irradiance, is the dominant influence on
where the MPP sits, even though irradiance dominates how much power is there.

\subsection{The dynamic resistance varies enormously}

At the MPP, $r_d$ ranges from \SI{10.4}{\ohm} at full sun to \SI{51.0}{\ohm} at
\SI{200}{\watt\per\square\metre}. Along the curve at fixed irradiance the swing
is far larger still, from under \SI{2}{\ohm} near $\Voc$ to nearly
\SI{500}{\ohm} on the current-source branch. Chapter~\ref{ch:converter} shows
that this variation maps directly onto the damping of the converter's dominant
mode, which is the central difficulty of the control problem.
